\documentclass[10pt, a4paper]{report}

\usepackage[utf8]{inputenc}
\usepackage[czech]{babel}
\usepackage{csquotes}
\usepackage[T1]{fontenc}
\usepackage{lmodern}

\usepackage{float}
\usepackage{caption}
\usepackage{nameref}
\usepackage{amsmath}
\usepackage{graphicx}
\usepackage[hidelinks]{hyperref}
\usepackage{titlesec}
\usepackage[
left=30mm,
right=30mm,
top=40mm,
bottom=30mm,
]{geometry}
\usepackage{todonotes}
\usepackage{array}
\usepackage{listings}
\usepackage{xcolor}

% Kapitoly bez velkého záhlaví "Kapitola N"
\titleformat{\chapter}{\normalfont\huge\bfseries}{\thechapter}{1em}{}
\titlespacing*{\chapter}{0pt}{0pt}{20pt}

% Kód (jen ASCII, listings s pdflatexem neumí UTF-8)
\lstset{
	basicstyle=\ttfamily\small,
	columns=fullflexible,
	keepspaces=true,
	frame=single,
	framerule=0.3pt,
	rulecolor=\color{gray},
	xleftmargin=2mm,
	xrightmargin=2mm,
	showstringspaces=false,
	breaklines=true,
}
\newcommand{\code}[1]{\texttt{#1}}

\begin{document}
	\sloppy % dlouhe nazvy v \code{} jinak pretekaji okraj
	
	\pagenumbering{arabic}
	\thispagestyle{empty}
	
	\begin{titlepage}
		\centering
		
		{\Large KIV/ZOS -- Základy operačních systémů} \\[0.5cm]
		{\Large Semestrální práce} \\[2.5cm]
		
		{\huge \bfseries PseudoFAT} \\[0.3cm]
		{\Large Zjednodušený souborový systém založený na FAT} \\[0cm]
		
		\vfill
		
		\begin{flushleft} \small
			Jméno: Jan Kopačka \\
			Osobní číslo: A24B0346P  \\
			Email: jankop@students.zcu.cz \\
			Datum: \today
		\end{flushleft}
		
	\end{titlepage}
	
	\setcounter{page}{2}
	
	\tableofcontents
	\newpage
	
	% =====================================================================
	\chapter{Úvod}
	
	\section{Účel projektu}
	
	Program implementuje zjednodušený souborový systém založený na FAT. Je napsaný v Rustu, což bylo pro tuto práci povoleno i přes požadavky zadání. Celý systém je
	uložený v jednom souboru na hostitelském disku a jeho obsah zůstává zachovaný mezi spuštěními.
	Program se spouští s jedním parametrem, cestou k tomuto souboru:
	
	\begin{lstlisting}
pseudofat myfs.dat
	\end{lstlisting}
	
	Ovládá se příkazy ze zadání (\code{format}, \code{mkdir}, \code{incp}, \code{cat} apod.)
	zadávanými na standardní vstup. Podle loginu (\code{jankop}) jsou navíc implementované příkazy
	\code{check} (kontrola konzistence) a \code{bug} (záměrné poškození). Úplné znění zadání je v repozitáři v~\code{docs/cz/Zadani.pdf}.
	
	Výstupy \code{OK}, textové výstupy (\code{ls}, \code{pwd}, \code{cat}, \code{info},
	\code{check}) a chyby parsování jdou na \code{stdout}. Chyby příkazů (\code{FILE NOT FOUND},
	\code{PATH NOT FOUND}, \code{EXIST}, \dots) jdou na \code{stderr}.
	
	Sestavení a spuštění zajišťuje \code{Makefile}: \code{make build} (release binárka do
	\code{dist/}), \code{make run}, \code{make test}, \code{make clean}.
	
	\section{Adresářová struktura}
	
	\begin{lstlisting}
Cargo.toml, Makefile
src/
  main.rs          vstupni bod, zpracovani argumentu
  shell.rs         smycka cteni prikazu, volani prikazu, vypis vysledku
  command.rs       parser prikazu (Command), Output, CommandExecError
  command/         jeden soubor = jeden prikaz (cat.rs, cp.rs, ..., load.rs)
  filesystem.rs    struktura FileSystem, FsError, operace se soubory
  filesystem/
    superblock.rs  superblok: rozlozeni disku, validace pri otevreni
    fat.rs         FAT tabulka: alokace, retezce, uvolnovani
    dir.rs         adresare a jejich polozky, jmena 8.3
    paths.rs       preklad cest
    consistency.rs logika pro prikazy check a bug
    disk.rs        nejnizsi vrstva: cteni/zapis bajtu na offsetu
  host.rs          cteni/zapis souboru na hostitelskem disku
  bytes.rs         serializace cisel (little-endian)
  config.rs        konstanty (velikost clusteru, limity)
  helpers.rs       pomocne funkce (varovani)
tests/
  integration.rs   integracni testy (spousti binarku, kontroluje vystup)
docs/cz/           tato dokumentace
	\end{lstlisting}
	
	\section{Architektura}
	
	Program je rozdělený do vrstev. Každá vrstva volá jen vrstvu pod sebou.
	
	\begin{table}[H]
		\centering
		\begin{tabular}{|l|l|p{6.2cm}|}
			\hline
			\textbf{vrstva} & \textbf{moduly} & \textbf{odpovědnost} \\ \hline
			shell & \code{shell}, \code{command} & čte řádky, parsuje je na \code{Command}, volá příkaz a vypisuje výsledek \\ \hline
			příkazy & \code{command/*} & tenké funkce \code{do\_*}: složí operace FS a převedou výsledek na \code{Output} \\ \hline
			souborový systém & \code{filesystem}, \code{filesystem/*} & superblok, FAT, adresáře, překlad cest, check a bug \\ \hline
			disk & \code{filesystem/disk} & čtení a zápis bajtů na daný offset souboru s FS \\ \hline
		\end{tabular}
	\end{table}
	
	Mimo tyto vrstvy stojí \code{host} (soubory na hostitelském disku pro \code{incp} a
	\code{outcp}), \code{bytes} (převod čísel na bajty) a \code{config} (konstanty).
	
	Struktura \code{FileSystem} drží otevřený disk, načtený superblok, celou FAT v paměti a cluster
	aktuálního adresáře. \code{Shell} drží \code{Option<FileSystem>}: \code{None} znamená
	nezformátovaný disk, kdy funguje jen \code{format} a ostatní příkazy skončí chybou
	\code{Disk is not formatted.}
	
	% =====================================================================
	\chapter{Architektonická rozhodnutí}
	
	\section{Návrh}
	
	\subsection{Rozložení disku}
	
	\begin{table}[H]
		\centering
		\begin{tabular}{|l|l|l|}
			\hline
			\textbf{offset} & \textbf{oblast} & \textbf{velikost} \\ \hline
			0 & superblok & 64 B \\ \hline
			\code{fat\_start} & FAT & $N \times 4$ B \\ \hline
			\code{data\_start} & clustery $0 \dots N-1$ (cluster 0 je kořen \code{/}) & $N \times 1024$ B \\ \hline
			& nevyužitý zbytek & méně než 1028 B \\ \hline
		\end{tabular}
	\end{table}
	
	Počet clusterů je $N = \lfloor (\text{disk\_size} - 64) / (1024 + 4) \rfloor$, protože každý
	cluster potřebuje 1024 B dat a 4 B ve FAT.
	
	Superblok (ve skutečném FAT nazývaný \enquote{boot sector}) popisuje celý disk. Program ho čte jako první
	a vše ostatní zjistí z něj. Všechna čísla jsou little-endian.
	
	\begin{table}[H]
		\centering
		\begin{tabular}{|r|r|l|l|p{4.5cm}|}
			\hline
			\textbf{offset} & \textbf{velikost} & \textbf{pole} & \textbf{příklad (600 MiB)} & \textbf{význam} \\ \hline
			0 & 8 B & \code{magic} & \code{PSEUDFAT} & poznání vlastního FS \\ \hline
			8 & 4 B & \code{version} & 1 & pro budoucí změny formátu \\ \hline
			12 & 4 B & \code{cluster\_size} & 1024 & \\ \hline
			16 & 4 B & \code{cluster\_count} & 612 009 & $N$ \\ \hline
			20 & 4 B & \code{root\_cluster} & 0 & cluster kořenového adresáře \\ \hline
			24 & 8 B & \code{disk\_size} & 629 145 600 & kontrola proti skutečné délce souboru \\ \hline
			32 & 8 B & \code{fat\_start} & 64 & začátek FAT \\ \hline
			40 & 8 B & \code{data\_start} & 2 448 100 & začátek clusteru 0 \\ \hline
			48 & 16 B & rezerva & nuly & \\ \hline
		\end{tabular}
	\end{table}
	
	Při otevření se ověří \code{magic} (jinak \enquote{nezformátováno}), \code{version},
	shoda \code{disk\_size} s délkou souboru a nakonec se rozložení spočítá znovu ze
	\code{disk\_size} a \code{cluster\_size} a musí se shodovat s uloženým. Poškozený superblok
	se tak nepoužije.
	
	\subsection{FAT tabulka}
	
	FAT je pole \code{cluster\_count} čísel \code{u32} hned za superblokem. Položka $i$ říká, co
	následuje po clusteru $i$:
	
	\begin{table}[H]
		\centering
		\begin{tabular}{|l|l|l|}
			\hline
			\textbf{hodnota na disku} & \textbf{\code{FATPtrEntry}} & \textbf{význam} \\ \hline
			\code{u32::MAX} & \code{FREE} & cluster je volný \\ \hline
			\code{u32::MAX - 1} & \code{EOF} & poslední cluster řetězce \\ \hline
			\code{u32::MAX - 2} & \code{BAD} & vadný cluster (jen pro úplnost, viz níže) \\ \hline
			cokoliv jiného & \code{Next(n)} & další cluster téhož souboru je $n$ \\ \hline
		\end{tabular}
	\end{table}
	
	\begin{itemize}
		\item FAT ví jen to, že po clusteru 2 následuje cluster 3. Kde soubor začíná, ve FAT není:
		první cluster je uložený v položce adresáře.
		\item Skutečný FAT čísluje datové clustery od 2, protože hodnota \code{FREE} je tam 0.
		Kdyby data mohla ležet v clusteru 0, znamenalo by \code{FAT[i] = 0} buď \enquote{volno},
		nebo \enquote{další je cluster 0}. Tady jsou speciální hodnoty na konci rozsahu \code{u32},
		takže se čísluje od 0. Z toho plyne i limit
		\code{MAX\_CLUSTER\_COUNT = u32::MAX - 2}.
		\item \code{BAD} je uvedená spíše pro úplnost. Zadání ji nevyžaduje a implementace ji plně
		nevyužívá: žádný cluster se jako vadný neoznačuje, jen \code{get\_chain} odmítne řetězec,
		který na \code{BAD} narazí.
		\item Skutečný FAT drží dvě kopie tabulky jako zálohu. Tady je jen jedna.
		\item Celá FAT je v paměti a každá změna se hned zapíše i na disk (write-through). Soubor
		s FS je tak po každé operaci aktuální.
	\end{itemize}
	
	\paragraph{Proč 32bitové položky.} FAT12 adresuje zhruba 4 080 clusterů, FAT16 zhruba
	65 520, FAT32 (28 bitů z 32) zhruba 268 milionů. Disk 600 MiB s clustery po 1 KiB má
	$600 \times 1024 = 614\,400$ clusterů (po odečtení režie zhruba 612 tisíc). To je víc, než
	zvládne FAT16, takže pro malé 1 KiB clustery připadají v úvahu jen 32bitové položky.
	Jedna FAT pak zabere $614\,400 \times 4\ \text{B} = 2\,457\,600\ \text{B} \approx 2{,}34$ MiB.
	Velikost FAT závisí na velikosti clusteru:
	
	\begin{table}[H]
		\centering
		\begin{tabular}{|r|r|r|r|}
			\hline
			\textbf{cluster} & \textbf{clusterů na 600 MiB} & \textbf{1 tabulka} & \textbf{2 tabulky} \\ \hline
			512 B & 1 228 800 & 4,69 MiB & 9,38 MiB \\ \hline
			1 KiB (použito) & 614 400 & 2,34 MiB & 4,69 MiB \\ \hline
			4 KiB (běžně FAT32) & 153 600 & 0,59 MiB & 1,17 MiB \\ \hline
		\end{tabular}
	\end{table}
	
	\subsection{Limity velikostí}
	
	\begin{itemize}
		\item \code{size} v položce adresáře je \code{u32}, jeden soubor má tedy nejvýš zhruba
		4 GiB. \code{cluster\_count} je také \code{u32}, ale omezuje celý disk
		(až $4{,}29 \cdot 10^9 \times 1024\ \text{B} \approx 4{,}4$ TB). Jde o dvě nezávislé
		veličiny se stejným stropem.
		\item Počet clusterů souboru se neukládá, dopočítá se jako
		$\lceil s / c \rceil$, kde $s$ je \code{size} a $c$ \code{cluster\_size}. Protože $c \ge 1$,
		platí $\lceil s / c \rceil \le s$: v nejhorším
		(cluster 1 B) jsou si rovny, jinak je clusterů méně. Když se \code{size} vejde do
		\code{u32}, vejde se tam i počet clusterů souboru.
		\item \code{MAX\_DISK\_SIZE = 4 GiB} je čistě vývojářský a testovací limit. Hlavně brání
		tomu, aby překlep ve \code{format} při testování vytvořil obrovský soubor. Na správnost nemá
		vliv: přetečení počtu clusterů hlídá \code{MAX\_CLUSTER\_COUNT} při vytváření superbloku
		a velikost souboru převod na \code{u32} v \code{create\_file}. Zadání požaduje 10--600 MB.
		\item \code{MIN\_CLUSTER\_SIZE = 512} B a podmínka mocniny dvou jsou pojistky spíš pro
		programátora, kdyby měnil \code{CLUSTER\_SIZE}. Používá se \code{CLUSTER\_SIZE = 1024} B.
	\end{itemize}
	
	\subsection{Adresáře}
	
	Adresář zabírá vždy právě jeden cluster (ve FAT má \code{EOF}, velikost 0) a nedá se zvětšovat.
	Cluster je rozdělený na sloty po 32 B. Každý slot je buď volný (samé nuly), nebo obsahuje
	položku:
	
	\begin{table}[H]
		\centering
		\begin{tabular}{|r|r|l|p{8cm}|}
			\hline
			\textbf{offset} & \textbf{velikost} & \textbf{pole} & \textbf{význam} \\ \hline
			0 & 12 B & \code{name} & 8 B jméno + 3 B přípona + \code{\textbackslash 0}, kratší části doplněné nulami (zadání: 12 B, formát 8.3) \\ \hline
			12 & 1 B & \code{kind} & 0 = volný slot, 1 = soubor, 2 = adresář (bajty 13--15 jsou nuly) \\ \hline
			16 & 4 B & \code{size} & velikost v bajtech; pro \code{info} a pro určení, kolik dat posledního clusteru je platných \\ \hline
			20 & 4 B & \code{start} & položka FAT: číslo prvního clusteru, nebo \code{EOF} (prázdný soubor) \\ \hline
			24 & 8 B & rezerva & \\ \hline
		\end{tabular}
	\end{table}
	
	Kapacita adresáře je $1024 / 32 - 2 = 30$ položek. Sloty 0 a 1 jsou \code{.} (odkaz na sebe) a
	\code{..} (odkaz na rodiče). Skutečný FAT je v kořeni nemá. Tady je má i kořen a jeho \code{..}
	ukazuje na něj samotného, takže \code{cd ..} v kořeni zůstane v kořeni. Uložení \code{.} a \code{..}
	stojí 2 sloty z 32, ale \code{cd ..} je pak jen čtení položky, \code{pwd} jde složit
	procházením nahoru a není potřeba držet cestu v paměti jako zásobník.
	
	\subsection{Překlad cesty}
	\label{sec:walk}
	
	Každý příkaz s cestou začíná stejně: z textu jako \code{../docs/./s1.txt} je potřeba najít
	položku.
	\begin{enumerate}
		\item Pokud cesta začíná \code{/}, začne se v kořeni, jinak v aktuálním adresáři.
		\item Cesta se rozdělí podle \code{/}, prázdné části (\code{a//b}, \code{a/}) se ignorují.
		\item Pro každou mezilehlou část: \code{.} znamená zůstat, \code{..} přejít do rodiče.
		Jakékoliv jiné jméno musí existovat a být adresář, jinak \code{PATH NOT FOUND}.
	\end{enumerate}
	
	Funkce \code{walk} projde smyčkou jen mezilehlé části a vrátí rodičovský adresář a poslední část
	jako \code{Option<\&str>}. Poslední část se odřízne, protože se na ni nesmí použít pravidlo
	\enquote{musí existovat a být adresář}: je to cíl operace a jeho pravidla určuje příkaz.
	
	\begin{table}[H]
		\centering
		\begin{tabular}{|l|l|l|}
			\hline
			\textbf{konec cesty} & \textbf{co to je} & \textbf{\code{walk} vrátí jako poslední část} \\ \hline
			\code{.../b}, \code{b} je adresář & existující adresář & \code{Some("b")} \\ \hline
			\code{.../f.txt}, \code{f.txt} je soubor & existující soubor & \code{Some("f.txt")} \\ \hline
			\code{.../new}, \code{new} neexistuje & nová položka (\code{mkdir}, \code{incp}) & \code{Some("new")} \\ \hline
			\code{.../.} nebo \code{.../..} & navigace & \code{None} (vyřešeno ve smyčce) \\ \hline
			\code{/} nebo prázdná cesta & kořen & \code{None} \\ \hline
		\end{tabular}
	\end{table}
	
	\code{walk} nerozhoduje mezi prvními třemi případy. Volající zpracuje poslední část jedním ze
	dvou způsobů:
	\begin{itemize}
		\item \textbf{najdi existující} (\code{cat}, \code{rm}, \code{cd}, \code{ls}, \code{info},
		zdroj u \code{cp}, \code{mv}, \code{outcp}): položka musí existovat a mít správný typ,
		\item \textbf{rozděl na rodiče a nové jméno} (\code{mkdir}, cíl u \code{cp}, \code{mv},
		\code{incp}): rodič musí existovat (jinak \code{PATH NOT FOUND}), jméno v něm ještě nesmí
		být (jinak \code{EXIST}) a musí splňovat formát 8.3.
	\end{itemize}
	
	\subsection{Pořadí zápisů a operace}
	
	Obecné pravidlo: při vytváření se zapíšou data, potom FAT a jako poslední položka adresáře. Při
	mazání se nejdřív odstraní položka a pak se uvolní FAT. Když program spadne uprostřed, nejhorší
	výsledek jsou obsazené clustery, na které nic neukazuje (ztracené clustery). To je jen ztráta
	místa. Při opačném pořadí by mohla zůstat položka ukazující na volné clustery, ty by pak přepsal
	jiný soubor a data dvou souborů by se smíchala.
	
	\paragraph{\code{incp s1 s2}} (\code{cp} funguje stejně, jen čte data z našeho FS):
	\begin{enumerate}
		\item Načte zdroj, jinak \code{FILE NOT FOUND}.
		\item Rozdělí cíl na rodiče a jméno (\code{PATH NOT FOUND}, \code{EXIST}, neplatné jméno).
		\item Ověří, že má rodič volný slot, jinak \code{Directory is full.}
		\item Ověří, že je volných $\lceil \text{size} / 1024 \rceil$ clusterů. Kontroluje se
		předem, aby nezůstal napůl zapsaný soubor.
		\item Zapíše data do volných clusterů (poslední bývá zaplněný jen částečně).
		\item Zřetězí je ve FAT: \code{FAT[c1] = c2}, \dots, \code{FAT[ck] = EOF}.
		\item Zapíše položku do rodiče.
	\end{enumerate}
	
	\paragraph{\code{rm s1}:} najde položku (musí to být soubor), odstraní ji z rodiče (vynuluje
	slot) a každý cluster řetězce nastaví na \code{FREE}. Data se při mazání nepřepisují, původní
	bajty v clusterech zůstanou, dokud je něco nepřepíše (na tom stojí nástroje pro obnovu
	smazaných souborů).
	
	\paragraph{\code{mkdir a1}:} rozdělí cestu na rodiče a jméno (stejné kontroly jako u
	\code{incp}), najde 1 volný cluster, nastaví \code{FAT[c] = EOF}, zapíše do clusteru \code{.}
	a \code{..} a nakonec položku do rodiče.
	
	\paragraph{\code{rmdir a1}:} najde adresář, který nesmí být kořen ani aktuální adresář (předek
	aktuálního adresáře neprojde, protože není prázdný). Kromě \code{.} a \code{..} musí být prázdný,
	jinak \code{NOT EMPTY}. Odstraní položku z rodiče a nastaví \code{FAT[c] = FREE}.
	
	\paragraph{\code{mv s1 s2}:} přesouvá se jen položka adresáře. Ze starého rodiče se odstraní a
	do nového se zapíše s novým jménem. Data ani FAT se nemění, takže přesun 500MB souboru trvá
	stejně dlouho jako přesun 1 B. Přejmenování je \code{mv} v rámci jednoho adresáře. U adresáře se
	navíc přepíše jeho \code{..} a je zakázaný přesun adresáře do sebe sama. \code{s2} je vždy
	úplná cílová cesta včetně jména, \code{mv a.txt dir} tedy při existujícím \code{dir} skončí
	\code{EXIST}.
	
	\subsection{check a bug}
	
	\code{check} projde strom adresářů od kořene a u každého souboru ověří neplatné číslo clusteru,
	cyklus v řetězci a nesoulad \code{size} s délkou řetězce. Každý problém vypíše na jeden řádek
	s cestou k souboru, bez problémů vypíše \code{OK}.
	
	\code{bug s1 [size|cycle|invalid]} poškodí soubor \code{s1}:
	
	\begin{table}[H]
		\centering
		\begin{tabular}{|l|p{5.5cm}|p{5cm}|}
			\hline
			\textbf{druh} & \textbf{co udělá} & \textbf{kdy jde použít} \\ \hline
			\code{size} & změní velikost v položce o jeden cluster & vždy \\ \hline
			\code{cycle} & poslední cluster nasměruje na první & soubor má aspoň 1 cluster a nepoškozený řetězec \\ \hline
			\code{invalid} & poslední cluster nasměruje na číslo mimo rozsah & soubor má aspoň 1 cluster a nepoškozený řetězec \\ \hline
		\end{tabular}
	\end{table}
	
	Bez druhého argumentu se druh vybere náhodně jen z druhů použitelných na daný soubor, takže
	opakované \code{bug s1} nikdy neselže kvůli nepoužitelné volbě. Náhoda bere nanosekundy
	systémového času, bez externí knihovny. \code{bug} lze volat opakovaně na různé soubory a
	\code{check} pak nahlásí všechna poškození.
	
	Poškozený soubor nejde smazat: \code{rm} napřed ověřuje řetězec, takže soubor s cyklem nebo
	neplatným clusterem skončí chybou. Jediná cesta zpět je \code{format}.
	
	\subsection{load}
	
	Skript pro \code{load} musí ležet v našem FS, nejdřív se tam tedy musí zkopírovat přes
	\code{incp}. Bez zformátovaného FS tak nefunguje nic kromě \code{format}. Chyba jednoho
	příkazu nezastaví zbytek skriptu. \code{exit} ve skriptu ukončí celý program. Vnořený
	\code{load} je zakázaný, protože skript načítající sám sebe by nikdy neskončil.
	
	\subsection{MiB a MB}
	
	Zadání uvádí velikost v MB (\code{format 600MB}). Program jednotku chápe jako mocninu dvou,
	tj. MiB: \code{600MB} vytvoří disk $600 \times 1024^2$ B. Správně by se tedy mělo psát MiB,
	zadání to tak ale nepíše a vyžaduje zápis \code{MB}. Proto se v textu i programu používá
	\code{MB}, přestože jde o MiB.
	
	\subsection{Multiplatformnost}
	
	Program funguje na Windows i Linuxu:
	\begin{itemize}
		\item všechna čísla na disku jsou natvrdo little-endian, soubor s FS je tak stejný na
		každé platformě (nezávisí na pořadí bajtů procesoru),
		\item argumenty a cesty se čtou jako \code{OsString}/\code{PathBuf} (Linux UTF-8,
		Windows UTF-16),
		\item konce řádků LF i CRLF se zpracují stejně (vstup i skripty pro \code{load}),
		\item výzva \code{pseudofat>} se vypisuje jen při čtení z terminálu, ne z roury,
		\item \code{Makefile} rozlišuje Windows a ostatní systémy.
	\end{itemize}
	
	% ---------------------------------------------------------------------
	\section{Implementace}
	
	\subsection{Příkazy jako \code{do\_*}}
	
	Každý příkaz je v samostatném souboru \code{command/<jméno>.rs} jako funkce \code{do\_<jméno>}
	v bloku \code{impl FileSystem}. Výjimky jsou dvě. \code{format} je volná funkce, protože pracuje i
	bez otevřeného FS. \code{load} je v bloku \code{impl Shell}, protože na rozdíl od ostatních
	příkazů vykonává další příkazy a víc spolupracuje se shellem (viz \ref{sec:pubcrate}). Logika, která potřebuje privátní pole \code{disk} a \code{fat}, je v
	\code{filesystem.rs} nebo jeho podmodulech (\code{paths.rs}, \code{consistency.rs}), protože
	podmodul vidí privátní pole rodiče. Příkazy pak jen skládají tyto operace.
	
	\subsection{Umělé \code{Next} v \code{get\_chain}}
	
	\code{get\_chain(first: FATPtrEntry)} nebere číslo clusteru, ale položku FAT. Pole
	\code{start} v položce adresáře je uložené jako \code{FATPtrEntry}: neprázdný soubor má
	\code{Next(první cluster)}, prázdný \code{EOF}. \code{get\_chain} pak se startem zachází jako
	s \enquote{ukazatelem na první cluster} a vrací \emph{všechny} clustery včetně prvního. Kdyby
	brala přímo první cluster, první by se vynechal nebo by se musel řešit zvlášť. Prázdný soubor
	dá bez zvláštní větve prázdný řetězec. Proto i \code{DirEntry::dir} zabaluje cluster adresáře do
	\code{Next}.
	
	\begin{lstlisting}
let clusters = self.fat.get_chain(entry.start)?;   
// start = Next(2) -> [2, 3, 4, 7, 10]
	\end{lstlisting}
	
	\subsection{\code{pub(crate)} v \code{shell.rs}}
	\label{sec:pubcrate}
	
	Z funkcí, které vykonávají příkazy (\code{run}, \code{execute}, \code{process\_line}), je
	\code{process\_line} jediná \code{pub(crate)}. Spolu s ní jsou \code{pub(crate)} i \code{Flow},
	\code{require\_fs}, \code{is\_load\_running} a \code{set\_load\_running}. Důvod: každý příkaz
	má být v samostatném souboru, ale \code{load} jako jediný potřebuje vnitřnosti \code{Shell}
	(vykonat řádek a vypsat výsledek). Jeho \code{impl Shell} v \code{command/load.rs} leží mimo
	modul \code{shell}, takže tyto položky musí být viditelné v rámci crate. Je to vědomá
	nekonzistence ve viditelnosti výměnou za jednotnou strukturu příkazů.
	
	\subsection{Dlouhé názvy v \code{paths.rs}}
	
	Funkce v \code{paths.rs} jsou velmi specifické a mají dlouhé ale dobře čitelné názvy
	(\code{get\_parent\_dir\_and\_name\_from\_path}, \code{get\_dir\_cluster\_from\_path},
	\code{get\_node\_from\_path}, \dots). Každá odpovídá jednomu způsobu zpracování poslední části
	cesty.
	
	\subsection{Další chytáky}
	
	\begin{itemize}
		\item \code{get\_node\_from\_path} vrací \code{Option}: \code{None} znamená kořen, který
		nemá rodiče ani položku se jménem. Adresář se hledá přes cluster (\code{dir\_of}), ne podle
		jména z cesty, protože cesta může končit \code{.} nebo \code{..} a pak skutečné jméno
		neobsahuje.
		\item \code{read\_file} odmítne soubor, jehož řetězec neodpovídá velikosti (jinak by vrátil
		nesmyslná data). \code{info} tuto kontrolu nedělá a nesoulad jen vypíše jako
		\code{MISMATCH}, aby šlo poškození po \code{bug} ukázat.
		\item V \code{move\_entry} jsou při přejmenování ve stejném adresáři \code{parent} a
		\code{target} dvě kopie téhož clusteru. Použije se jen \code{replace} na jedné z nich
		(funguje i v plném adresáři). Zápis přes obě kopie by jednu změnu přepsal druhou.
		\item Smyčky, které stoupají přes \code{..} (\code{is\_dir\_inside},
		\code{get\_path\_from\_cluster}) a prochází řetězec (\code{get\_chain}), mají strop
		\code{cluster\_count} kroků. V poškozeném FS by jinak cyklus znamenal nekonečnou smyčku.
	\end{itemize}
	
	% =====================================================================
	\chapter{Největší problémy}
	
	\paragraph{Zabránění \code{load} uvnitř \code{load}.} Skript, který načítá sám sebe, by běžel
	donekonečna. \code{Shell} má příznak \code{loading}, který se nastaví po dobu běhu skriptu.
	Vnořený \code{load} pak skončí chybou \code{nested load}.
	
	\paragraph{Ukládání změn FAT (RAM a disk).} FAT je celá v paměti kvůli rychlému hledání
	volných clusterů a procházení řetězců. Každá změna se ale hned zapíše i na disk
	(write-through), aby soubor s FS nikdy nezaostal za stavem v paměti. Zapisuje se jen změněný
	úsek tabulky, ne celá FAT.
	
	\paragraph{Překlad cesty na cluster.} Najít pro \code{/docs/images} cluster adresáře znamená
	projít části cesty od kořene nebo aktuálního adresáře, řešit \code{.} a \code{..} a poslední část
	nechat na volajícím (viz \nameref{sec:walk}). Hlavní potíž byla oddělit mezilehlé části (musí
	být adresáře) od poslední části (pravidla určuje příkaz).
	
	\paragraph{Cesta ke clusteru (\code{pwd}).} Adresář nezná své jméno, to je uložené v rodiči.
	Cesta se proto skládá odspodu: přes \code{..} se přejde do rodiče, v něm se najde položka, která
	ukazuje na výchozí cluster, a tak dál až ke kořeni. Jména se sbírají od nejhlubšího a na konci se
	otočí.
	
	\paragraph{Efektivní zápis souboru.} Clustery souboru se seskupí do souvislých úseků (extentů) a
	každý úsek se zapíše nebo přečte jedním voláním. Soubor v clusterech \code{[4, 5, 9]} se tak
	zapíše dvěma zápisy místo tří. Drobnost: změny FAT se ukládají jako jeden úsek od nejmenšího po
	největší dotčený cluster. Když má soubor jeden cluster na začátku a druhý na konci disku,
	přepíše se celá FAT.
	
	\paragraph{Přesun adresáře.} Při \code{mv} adresáře nestačí přesunout položku. Musí se přepsat
	jeho \code{..} na nového rodiče (jinak by \code{cd ..} a \code{pwd} vedly do starého
	umístění), zakázat přesun do vlastního podstromu (vznikla by smyčka odpojená od kořene) a
	všechny kontroly provést dřív, než se cokoliv zapíše.
	
\end{document}
